\uuid{KG6y}
\titre{Image réciproque}
\auteur{bodin ; hanani ; mzari}
\organisation{exo7}
\licence{CC BY-NC-SA 4.0}
\datecreate{2020-01-20}
\niveau{L1}
\module{Algèbre}
\chapitre{Logique, ensemble, raisonnement}
\sousChapitre{Ensemble}
\difficulte{2}
\theme{Ensembles, applications, image réciproque}
\exo7qcmid{lille-1.ensembles.017}

\contenu{
\begin{question}
\qclassification{100.02,101.01,102.02,102.02}
\qtags{motcle=image réciproque}
On considère l'application $f:\Rr\to \Rr$ définie par
$$\forall x\in \Rr,\; f(x)=x^2+1.$$
Quelles sont les bonnes réponses ?
\begin{answers}
    \good{$f^{-1}([1,5])=[-2,2]$}
    \good{$f^{-1}([0,5])=[-2,2]$}
    \bad{$f^{-1}([1,5])=[0,2]$}
    \bad{$f^{-1}([0,5])=[0,2]$}
\end{answers}
\begin{explanations}
D'une part, $x\in f^{-1}([1,5])\Leftrightarrow f(x)\in [1,5]\Leftrightarrow x^2\leq 4$. D'autre part, $x\in f^{-1}([0,5])\Leftrightarrow f(x)\in [0,5]\Leftrightarrow x^2\leq 4$. Donc $f^{-1}([1,5])=f^{-1}([0,5])=[-2,2]$.
\end{explanations}
\end{question}
}
